\paragraph{Separable extensions of a one-owner trace.}
The interval criterion also has a closed form for one-owner bases,
without an acyclicity assumption. Let $V$ be finite and nonempty and
let $c$ be a loopless closure on $V$. Fix an owner $r$. Define $C$ as
the union of the $\star_r$-folds over the sets $V\setminus H$, for
all proper $c$-closed sets $H$. These sets cover $V$, since
$c(\varnothing)=\varnothing$. This one-owner trace has
\begin{equation}\label{core:eq:one-owner-trace-value}
 r\in C(w)\iff r\in\bigcup_iw_i,\qquad
 k\in C(w)\iff
 \exists i\notin c(E_r(w))\quad k\in w_i\quad(k\ne r).
\end{equation}
Indeed an $r$-fold preserves $r$ and supplies a foreign winner precisely
when its hyperedge avoids $E_r(w)$. A closed set containing $E_r(w)$
and missing $i$ exists exactly when $i\notin c(E_r(w))$.

Take normalized antitone factors $B_s:2^V\to\{0,1\}$, with
$B_s(\varnothing)=1$, and form the source-reconstructing candidate
\[
 F(X,w)=C(w)\cup
 \{k\in X:\bigwedge_{s\ne k}B_s(E_s(w))=1\}.
\]
The passive source predicates of $F$ equal those of $C$.
Colored-cofactor independence means the law in
\cref{core:cor:colored-cofactor-dependence}: whenever a source predicate
vanishes at a proper partial atomic coloring, the corresponding
cofactor is independent of that source's entire menu.

\begin{theorem}[Separable intervals above one-owner traces]
\label{core:thm:one-owner-source-interval}
For $K\geq3$, the candidate $F$ is core if and only if it satisfies
colored-cofactor independence. Equivalently, writing $\psi_s=\neg B_s$,
the following two subset laws are necessary and sufficient:
\begin{align*}
 \psi_r(c(E))&=\psi_r(E)&&\text{for all }E\subseteq V,\\
 \psi_s(A\setminus c(E))&=\psi_s(A)
 &&\text{if }s\ne r,\ E\cap A=\varnothing,\ \psi_r(E)=0.
\end{align*}
In that case there is a primitive spine
$S$ with exactly one occurrence of $X$, full type $(B_s)_s$, and
$S\subseteq F$, such that $F=C+S$.
\end{theorem}
\begin{proof}
Core realizability implies colored-cofactor independence. We first
deduce the two subset laws from independence and then construct a
realizer from those laws alone. Each nonzero $\psi_s$ is a normalized monotone
Boolean function. We omit the helper for every identically zero
$\psi_s$.

Fix a foreign queried winner $k$, set
$X=\{k\}$, fix $E$ to $r$, and make the remaining inputs neutral
for $k$. The baseline omits $k$, and its captured inputs are
$c(E)\setminus E$. Independence lets us recolor all those inputs
to $r$ without changing the answer. Hence
\[
 \psi_r(c(E))=\psi_r(E).
\]
For the second law choose $k\notin\{r,s\}$, put $X=\{k\}$,
color $E$ by $r$ and $A$ by $s$, and make the other passive inputs
neutral for $k$. The $k$-answer is $B_s(A)$, since $B_r(E)=1$.
Every input in $A\cap c(E)$ is captured after fixing $E$ to $r$.
Independence lets us neutralize them all without changing the answer,
which proves the second law.

Now assume the two laws. When $\psi_r\ne0$, the first shows that
$c(E)=V$ implies $\psi_r(E)=1$. In this case define
\[
 c_f(E)=\begin{cases}V,&\psi_r(E)=1,\\c(E),&\psi_r(E)=0\end{cases}
\]
This is a loopless closure extending $c$, with spanning predicate $\psi_r$.
Its one-owner trace $H_r$ is bounded by $C$, by
\eqref{core:eq:one-owner-trace-value}, and has
$\Phi_r(H_r)=\psi_r$.

For $s\ne r$ with $\psi_s\ne0$, let $d_s$ be the Horn closure
generated by the implications
\begin{equation}\label{core:eq:conditional-horn-helpers}
 E\longrightarrow c(E)\quad\text{if }\psi_r(E)=0,
 \qquad
 E\longrightarrow V\quad\text{if }\psi_s(V\setminus E)=0.
\end{equation}
Here $\psi_r=0$ is allowed. We claim
\begin{equation}\label{core:eq:conditional-horn-span}
 d_s(T)=V\quad\Longleftrightarrow\quad\psi_s(V\setminus T)=0.
\end{equation}
The reverse implication is a generating rule. For the forward
implication assume $\psi_s(V\setminus T)=1$ and follow closure
generation. A collapse rule cannot apply to any $E\subseteq T$,
by monotonicity. Consider instead $E\subseteq T$ with $\psi_r(E)=0$.
The second subset law, with $A=V\setminus T$, gives
$\psi_s(A\setminus c(E))=1$.
Thus adding $c(E)$ preserves the same invariant.
Thus generation cannot reach $V$, where the invariant would read
$\psi_s(\varnothing)=1$. This proves
\eqref{core:eq:conditional-horn-span}. It also shows that $d_s$ is
loopless: no collapse rule applies at $\varnothing$, and
$c(\varnothing)=\varnothing$.

For every proper $d_s$-closed set $H$, take the $\star_r$-fold
over $V\setminus H$, and combine these clauses by $\star_s$.
Call the result $H_s$. Use all proper closed sets here; maximal
closed sets need not determine an arbitrary closure.
The inner clauses and outer fold give
\[
 \Phi_s(H_s)(A)=[d_s(V\setminus A)\ne V]=\psi_s(A).
\]
They also show that any foreign-$r$ winner supplied by position $i$
requires $i\notin d_s(E_r(w))$. If $\psi_r(E_r(w))=0$, the first
rules in \eqref{core:eq:conditional-horn-helpers} give
$d_s(E_r(w))\supseteq c(E_r(w))$, so every such winner lies in
$C(w)$. Any $r$ supplied by $H_s$ is already in some input and
hence in $C(w)$. Consequently
\[
 H_s(w)\subseteq C(w)\quad\text{whenever }\psi_r(E_r(w))=0.
\]

Start a spine at $X$, attach these $H_s$ by $\star_s$ in any
fixed order, and attach $H_r$ last by $\star_r$ if present.
Its marked type is exactly $(B_s)_s$.
On the stratum $\psi_r(E_r)=0$, all helpers are bounded by $C$,
so the spine lies in $C\cup X$. On the other stratum the final
$H_r$ equals $\{r\}$ and absorbs the prefix; this winner belongs
to $C$. Thus $S\subseteq C\cup X$ everywhere. Union with $C$
has the prescribed distinguished source predicate and the baseline's
passive predicates, hence equals $F$. In particular $S\subseteq F$.
\end{proof}

\paragraph{The terminal interface over several owners.}
The functional-independence part has a closed form for any
source-separable core baseline $C$, with owner closures $(c_t)_t$.
Adding $X$ as above leaves every passive source predicate unchanged;
thus $X$ is a source sink. Write $\psi_t=\neg B_t$.

\begin{lemma}[Terminal cofactor independence]
\label{core:lem:terminal-closure-laws}
The candidate satisfies colored-cofactor independence if and only if
\begin{align*}
 \psi_t(c_t(E))&=\psi_t(E),\\
 \psi_s(A\setminus c_t(E))&=\psi_s(A)
 \quad(s\ne t,\ E\cap A=\varnothing,\ \psi_t(E)=0).
\end{align*}
Under these laws its full value is
\[
 F(X,w)=\begin{cases}
 C(w)\cup X,&\psi_t(E_t(w))=0\text{ for every }t,\\
 C(w),&\text{otherwise}.
 \end{cases}
\]
\end{lemma}
\begin{proof}
The pure-$t$ probe and the two-color $(t,s)$ probe in the preceding
proof establish necessity, with $c_t$ in place of $c$.

For sufficiency put $M(w)=\{t:\psi_t(E_t(w))=1\}$.
If $M=\{t\}$, then $t\in C(w)$. Otherwise every $j\in E_t(w)$
is captured by some foreign $c_s(E_s(w))$, since $C$ is separable
and omits $t$. The second law, using $\psi_s(E_s)=0$, lets us
remove those captured inputs from $E_t$ while preserving
$\psi_t(E_t)=1$. All of $E_t$ is removed, contradicting
$\psi_t(\varnothing)=0$.
When $|M|\geq2$, no winner passes all its foreign-owner masks;
when $M=\{t\}$, only $t$ can pass, and it is already in $C$.
This proves the value formula.

If the new-source predicate vanishes at a proper atomic fixing, some
mask is already positive and stays so under every extension of that
fixing. The formula makes the whole cofactor equal to $C$, independently
of $X$.

For a captured passive source $j$, fix the other menus and first give
$j$ an ambiguous menu. Some owner $t$ captures it, so
$j\in c_t(E_t)$. The baseline is independent of $j$ under this fixing.
If $\psi_t(E_t)=1$, a mask stays positive when $j$ is changed to any
atomic menu, so $F=C$ throughout. Otherwise the first law makes
$\psi_t$ invariant under adding $j$ to $E_t$. For every $s\ne t$,
apply the second law to both $A=E_s$ and $A=E_s\cup\{j\}$;
their differences with $c_t(E_t)$ coincide, so $\psi_s$ is invariant
too. Hence $M$ is unchanged by any menu of $j$, and the displayed
formula proves independence. This covers every captured passive
cofactor and completes the proof.
\end{proof}

The lemma characterizes the functional interface. The realization
theorem above additionally uses that the baseline is a one-owner trace.
